% Generated draft: requires mathematical and source review.
\documentclass[11pt]{article}
\usepackage[margin=25mm]{geometry}
\usepackage{fontspec}
\setmainfont{Latin Modern Roman}
\usepackage{amsmath,amssymb,mathtools,unicode-math}
\setmathfont{Latin Modern Math}
\usepackage{enumitem,needspace,xurl,hyperref}
\setlength{\parindent}{0pt}
\setlength{\parskip}{5pt}
\setlength{\emergencystretch}{3em}
\newcommand{\N}{\mathbb{N}}
\newcommand{\Z}{\mathbb{Z}}
\newcommand{\Q}{\mathbb{Q}}
\newcommand{\R}{\mathbb{R}}
\newcommand{\C}{\mathbb{C}}
\newcommand{\F}{\mathbb{F}}
\newcommand{\A}{\mathbb{A}}
\newcommand{\PP}{\mathbb{P}}
\newcommand{\E}{\mathbb{E}}
\newcommand{\eps}{\varepsilon}
\newcommand{\dd}{\,\mathrm{d}}
\begin{document}
\section*{UnsolvedMath ID: 701044}
\textbf{Source problem code:} catalogue entry
\subsection{Definitions and mathematical statement}
Let \(G\) act on a set \(\Omega\) with \(|\Omega|\geq 3\); no finiteness assumption is imposed. For distinct points \(\alpha_1,\ldots,\alpha_k\in\Omega\), write
\[
G_{\alpha_1,\ldots,\alpha_k}=\{g\in G:g\alpha_i=\alpha_i\text{ for every }i\}.
\]
Here stabilizers are interpreted as pointwise stabilizers. Assume that the action is triply transitive: for any two ordered triples of pairwise distinct points, there is an element of \(G\) carrying the first triple to the second. Assume also that, for every pair of distinct points \(\alpha,\beta\),
\[
\{g\in G_{\alpha,\beta}:g^2=1\}=\{1\},
\]
and that \(G_{\alpha,\beta,\gamma}=\{1\}\) for every triple of pairwise distinct points. Thus the action is sharply triply transitive and faithful.

For a field \(P\), let
\[
\operatorname{PGL}_2(P)=\operatorname{GL}_2(P)/\{\lambda I_2:\lambda\in P^\times\}.
\]
Its natural action on \(\mathbb P^1(P)=P\cup\{\infty\}\) is the action on one-dimensional subspaces of \(P^2\), identifying \(x\in P\) with \([x:1]\) and \(\infty\) with \([1:0]\).

\emph{Question (V. D. Mazurov, Kourovka Notebook Problem 14.59).} Must there exist a field \(P\) of characteristic \(2\), a group isomorphism \(\varphi:G\to\operatorname{PGL}_2(P)\), and a bijection \(f:\Omega\to\mathbb P^1(P)\) such that
\[
f(g\alpha)=\varphi(g)f(\alpha)
\qquad(g\in G,\ \alpha\in\Omega)?
\]
This interprets the source's word ``similar'' as isomorphism of permutation actions. The source does not explicitly specify these stabilizer and similarity conventions, so they require confirmation against the original text.

The supplied statement reports an affirmative answer under the additional hypothesis that a two-point stabilizer is periodic, meaning that every element of that stabilizer has finite order. No uniform bound on these orders is stipulated. This special-case assertion is retained as supplied source information, without independent verification.
\subsection{Short English statement}
Must every sharply triply transitive permutation group whose pointwise two-point stabilizers contain no involutions be permutation-isomorphic to the natural action of \(\operatorname{PGL}_2(P)\) for some field \(P\) of characteristic \(2\)? The supplied source reports that this holds when a two-point stabilizer is periodic.
\subsection{Sources}
\begin{itemize}
\item Local UnsolvedMath registry, permanent ID \(701044\), exact source problem code \texttt{catalogue entry}, title ``Kourovka Notebook Problem 14.59.'' The supplied registry labels the problem \texttt{open}; that label is not an independently verified literature assessment.
\item \emph{The Kourovka Notebook (Unsolved Problems in Group Theory)}, Problem 14.59, proposed by V. D. Mazurov, 14th issue (1999). The supplied background identifies the referenced version as the 21st issue edition (2026), \texttt{arXiv:1401.0300v48}. The supplied collection description attributes editorship to E. I. Khukhro and V. D. Mazurov. This bibliographic information is reproduced from the record; the source has not been independently consulted.
\item The supplied problem statement reports the affirmative result for periodic two-point stabilizers, but supplies no separate citation or attribution for that result.
\end{itemize}
\end{document}
